% GENERATED FILE. Edit canonical lessons, then run npm run generate:books.
% canonical-source-hash: fnv1a64:7c89bd36a74488a8
% canonical-lessons: FR-C163-appartement, FR-C163-salle-a-manger, FR-C163-balcon, FR-C163-serrure, FR-C163-sonnette

\chapter{Flat, Dining room, Balcony, Lock, Doorbell}
\label{ch:fr-flat-dining-room-balcony-lock-doorbell}
\hlchaptermodality{163}

\begin{chapteropening}
\textbf{By the end of this chapter:} \emph{I can say a flat, a dining room, a balcony, a lock and a doorbell.}

Complete the last lesson of chapter 163: I can say a flat, a dining room, a balcony, a lock and a doorbell.
\end{chapteropening}

\section[l'appartement]{l'appartement — a flat}

\label{lesson:FR-C163-appartement}



\begin{quote}
Before the new one: say the French for married, then the French for divorced.
\end{quote}

\subsection*{You'll want to know: l'appartement}
\textbf{l'appartement} — \textquotedblleft{}a flat\textquotedblright{}.

\begin{cousinweb}
\textbf{appartement} $\leftarrow$ Italian \emph{appartamento} (EtymDB 2.1) — borrowed; built on Latin \emph{ad} and \emph{pars}, \textquotedblleft{}a part\textquotedblright{}: a set of rooms apart.
\end{cousinweb}

\begin{grammarlens}[title={What you've built}]
One more word for this chapter.
\end{grammarlens}

\subsection*{Your turn}
\begin{itemize}
\raggedright
  \item \emph{Say it:} \emph{l'appartement}
  \item \emph{Say it:} \emph{l'appartement}, once more
  \item \emph{Say it:} say \emph{l'appartement}
  \item \emph{Recall:} say the French for married, then the French for divorced, then say \emph{l'appartement} again
\end{itemize}

\begin{tcolorbox}[breakable,colback=teal!4,colframe=teal!35!black,title={Before you move on}]
What does l'appartement mean? (\textquotedblleft{}A flat\textquotedblright{}.) Say it once more.
\end{tcolorbox}
\section[la salle à manger]{la salle à manger — a dining room}

\label{lesson:FR-C163-salle-a-manger}



\begin{quote}
Before the new one: say the French for divorced, then the French for a flat.
\end{quote}

\subsection*{You'll want to know: la salle à manger}
\textbf{la salle à manger} — \textquotedblleft{}a dining room\textquotedblright{}.

\begin{cousinweb}
\textbf{salle à manger} — \textquotedblleft{}room for eating\textquotedblright{}: \emph{salle}, \textquotedblleft{}a room\textquotedblright{}, \emph{à}, \textquotedblleft{}for\textquotedblright{}, and \emph{manger}, \textquotedblleft{}to eat\textquotedblright{}. \emph{Salle} $\leftarrow$ Frankish \emph{sal}, \textquotedblleft{}hall\textquotedblright{}, \emph{à} $\leftarrow$ Latin \emph{ad}, and \emph{manger} $\leftarrow$ Latin \emph{manducāre}, \textquotedblleft{}to chew\textquotedblright{} (EtymDB 2.1).
\end{cousinweb}

\begin{grammarlens}[title={What you've built}]
One more word for this chapter.
\end{grammarlens}

\subsection*{Your turn}
\begin{itemize}
\raggedright
  \item \emph{Say it:} \emph{la salle à manger}
  \item \emph{Say it:} \emph{la salle à manger}, once more
  \item \emph{Say it:} say \emph{la salle à manger}
  \item \emph{Recall:} say the French for divorced, then the French for a flat, then say \emph{la salle à manger} again
\end{itemize}

\begin{tcolorbox}[breakable,colback=teal!4,colframe=teal!35!black,title={Before you move on}]
What does la salle à manger mean? (\textquotedblleft{}A dining room\textquotedblright{}.) Say it once more.
\end{tcolorbox}
\section[le balcon]{le balcon — a balcony}

\label{lesson:FR-C163-balcon}



\begin{quote}
Before the new one: say the French for a flat, then the French for a dining room.
\end{quote}

\subsection*{You'll want to know: le balcon}
\textbf{le balcon} — \textquotedblleft{}a balcony\textquotedblright{}.

\begin{cousinweb}
\textbf{balcon} $\leftarrow$ Italian \emph{balcone}, from Lombardic \emph{*balko}, \textquotedblleft{}a beam\textquotedblright{} — borrowed. (EtymDB 2.1)
\end{cousinweb}

\begin{grammarlens}[title={What you've built}]
One more word for this chapter.
\end{grammarlens}

\subsection*{Your turn}
\begin{itemize}
\raggedright
  \item \emph{Say it:} \emph{le balcon}
  \item \emph{Say it:} \emph{le balcon}, once more
  \item \emph{Say it:} say \emph{le balcon}
  \item \emph{Recall:} say the French for a flat, then the French for a dining room, then say \emph{le balcon} again
\end{itemize}

\begin{tcolorbox}[breakable,colback=teal!4,colframe=teal!35!black,title={Before you move on}]
What does le balcon mean? (\textquotedblleft{}A balcony\textquotedblright{}.) Say it once more.
\end{tcolorbox}
\section[la serrure]{la serrure — a lock}

\label{lesson:FR-C163-serrure}



\begin{quote}
Before the new one: say the French for a dining room, then the French for a balcony.
\end{quote}

\subsection*{You'll want to know: la serrure}
\textbf{la serrure} — \textquotedblleft{}a lock\textquotedblright{}.

\begin{cousinweb}
\textbf{serrure} — built from the verb for \textquotedblleft{}to squeeze, to shut\textquotedblright{}, $\leftarrow$ Late Latin \emph{serare}, \textquotedblleft{}to bolt\textquotedblright{}, from \emph{sera}, \textquotedblleft{}a door bar\textquotedblright{} (Lewis \& Short).
\end{cousinweb}

\begin{grammarlens}[title={What you've built}]
One more word for this chapter.
\end{grammarlens}

\subsection*{Your turn}
\begin{itemize}
\raggedright
  \item \emph{Say it:} \emph{la serrure}
  \item \emph{Say it:} \emph{la serrure}, once more
  \item \emph{Say it:} say \emph{la serrure}
  \item \emph{Recall:} say the French for a dining room, then the French for a balcony, then say \emph{la serrure} again
\end{itemize}

\begin{tcolorbox}[breakable,colback=teal!4,colframe=teal!35!black,title={Before you move on}]
What does la serrure mean? (\textquotedblleft{}A lock\textquotedblright{}.) Say it once more.
\end{tcolorbox}
\section[la sonnette]{la sonnette — a doorbell}

\label{lesson:FR-C163-sonnette}



\begin{quote}
Before the new one: say the French for a balcony, then the French for a lock.
\end{quote}

\subsection*{You'll want to know: la sonnette}
\textbf{la sonnette} — \textquotedblleft{}a doorbell\textquotedblright{}.

\begin{cousinweb}
\textbf{sonnette} — built from the verb for \textquotedblleft{}to ring\textquotedblright{}, $\leftarrow$ Latin \emph{sonare}, \textquotedblleft{}to sound\textquotedblright{}, with \emph{-ette}, \textquotedblleft{}little\textquotedblright{}. (EtymDB 2.1)
\end{cousinweb}

\begin{grammarlens}[title={What you've built}]
One more word for this chapter.
\end{grammarlens}

\subsection*{Your turn}
\begin{itemize}
\raggedright
  \item \emph{Say it:} \emph{la sonnette}
  \item \emph{Say it:} \emph{la sonnette}, once more
  \item \emph{Say it:} say \emph{la sonnette}
  \item \emph{Recall:} say the French for a balcony, then the French for a lock, then say \emph{la sonnette} again
\end{itemize}

\begin{tcolorbox}[breakable,colback=teal!4,colframe=teal!35!black,title={Before you move on}]
What does la sonnette mean? (\textquotedblleft{}A doorbell\textquotedblright{}.) Say it once more.
\end{tcolorbox}
